\subsection{Concentration and Stability}

Generalization can also follow from concentration of the training criterion or from stability of the learning algorithm itself.

\begin{definition} \label{def:sl-subgaussian-random-variable} \hlt{Subgaussian Random Variable}\\
A centered random variable $Z$ is \hlt{subgaussian} with variance proxy $\sigma^2$ if
\[
\mathbb{E}\exp(\lambda Z)\leq \exp(\lambda^2\sigma^2/2)
\]
for every $\lambda\in\R$.  Subgaussian variables have Gaussian-type tail bounds.
\end{definition}

\begin{theorem} \label{thm:sl-mcdiarmid-form} \hlt{Bounded Differences Inequality}\\
Let $F(X_1,\ldots,X_n)$ change by at most $c_i$ when only $X_i$ is replaced.  Then
\[
P\{F-\mathbb{E}F\geq t\}
\leq
\exp\left(-\frac{2t^2}{\sum_{i=1}^nc_i^2}\right).
\]
\end{theorem}

\begin{definition} \label{def:sl-uniform-stability} \hlt{Uniform Stability}\\
A learning algorithm $A$ has \hlt{uniform stability} $\beta_n$ for loss $\ell$ if replacing one training example changes the loss on any test point by at most $\beta_n$:
\[
\sup_z\abs{\ell(A(S),z)-\ell(A(S'),z)}\leq\beta_n.
\]
\end{definition}

\begin{remark} \label{rem:sl-stability-generalization} \hlt{Algorithm-Dependent Generalization}\\
Stability bounds depend on the training procedure, not just the hypothesis class.  This makes them useful for regularized empirical risk minimization and for understanding why some optimization algorithms generalize better than the raw size of their model class suggests.
\end{remark}
